\documentclass[11pt]{article}
\usepackage[margin=1.2in]{geometry}
\usepackage{amsmath,amssymb,amsthm}
\usepackage{mathrsfs}
\usepackage{booktabs}
\usepackage[colorlinks=true,linkcolor=blue,citecolor=blue]{hyperref}

\newtheorem{theorem}{Theorem}[section]
\newtheorem{proposition}[theorem]{Proposition}
\newtheorem{lemma}[theorem]{Lemma}
\newtheorem{condition}[theorem]{Condition}
\newtheorem{claim}[theorem]{Claim}
\newtheorem{remark}[theorem]{Remark}

\newcommand{\R}{\mathbb{R}}
\newcommand{\sG}{\mathscr{G}}
\newcommand{\sS}{\mathscr{S}}
\newcommand{\bN}{\mathbf{N}}
\newcommand{\bM}{\mathbf{M}}
\newcommand{\E}{\mathbb{E}}
\newcommand{\PP}{\mathbb{P}}
\newcommand{\cE}{\mathcal{E}}
\newcommand{\dd}{\,\mathrm{d}}
\DeclareMathOperator{\ent}{ent}

\title{The storage capacity of the Ising perceptron:\\
verification of the outstanding numerical conditions}
\author{Yitzchak Shmalo\\
\small Einstein Institute of Mathematics, The Hebrew University of Jerusalem\\
\small yitzchak.shmalo@gmail.com}
\date{July 2026}

\begin{document}
\maketitle

\begin{abstract}
Krauth and M\'ezard (1989) predicted that the storage capacity of the Ising
perceptron is $\alpha_\star N + o(N)$ for an explicit constant
$\alpha_\star \approx 0.833$. Ding and Sun proved the matching lower bound,
and Huang recently proved the matching upper bound; each result is
conditional on a numerical hypothesis about an explicit low-dimensional
variational problem, stated by those authors and checked by them only in
floating point. Huang's own rigorous interval arithmetic verified every
numerical constant in his upper bound \emph{except} the main
Condition~1.3---the nonpositivity, over the whole plane, of a two-variable
rate function whose supremum is not attained. That unbounded, degenerate
optimization is the reason the capacity upper bound remained conditional.

We verify the condition in full. Reparametrizing Huang's functional into the
moment coordinates of the pair $(X,\tanh X)$ compresses the plane onto a
\emph{compact} convex body and turns the entropy term into a per-cell linear
function by convex duality; certified sweeps then establish nonpositivity on
the whole body away from a star-shaped region around the single degenerate
maximizer. On the star itself no box enclosure can work---the tilted
covariance $\nabla^2\Phi$ is nearly singular and its inverse loses an
essential cancellation---so we certify concavity of a majorant along rays
from the maximizer, bounding the entropy Hessian by the explicit matrix
$\E[ff^\top\max_\Lambda \mathrm{sech}^2(\lambda\cdot f)]$ over sublevel-set
localizations of the dual that are themselves certified by one-dimensional
sweeps; nothing is ever inverted. The same machinery, in one variable,
verifies the Ding--Sun Condition~1.2 including the degenerate zero, and we
re-establish the parameter rectangle
$(\alpha_\star, q_\star, \psi_\star)$ on which both papers rest. Together
these make the Krauth--M\'ezard prediction for the Ising perceptron capacity
unconditional.
\end{abstract}

\section{Introduction}

Fix $\kappa = 0$ and let $g^1, g^2, \ldots$ be independent standard Gaussian
vectors in $\R^N$. The Ising perceptron is the random subset of the discrete
cube
\[
  S_N^M \;=\; \bigl\{ x \in \{-1,+1\}^N : \langle g^a, x\rangle \ge 0
  \text{ for all } 1 \le a \le M \bigr\},
\]
and its capacity is $M_N/N$, where $M_N$ is the largest $M$ with
$S_N^M \neq \emptyset$. Krauth and M\'ezard \cite{krauth1989storage} analyzed
this model with the replica method and predicted that $M_N/N$ converges in
probability to an explicit constant $\alpha_\star \approx 0.83307860$,
defined below. The prediction has been one of the standard test problems for
making the cavity method rigorous ever since.

Rigorous progress converged on the constant from both sides. Ding and
Sun \cite{ding2018capacity} proved
$\liminf_N \PP(M_N/N \ge \alpha) > 0$ for every $\alpha < \alpha_\star$;
combined with the sharp threshold sequence of Xu \cite{xu2021sharp} and
Nakajima and Sun \cite{nakajima2023sharp}, this gives the lower bound with
high probability. Huang \cite{huang2024capacity} proved
$\PP(M_N/N \ge \alpha) \to 0$ for every $\alpha > \alpha_\star$. Both
results are conditional. The Ding--Sun theorem assumes their Condition~1.2:
a certain function $\lambda \mapsto \mathscr{S}_\star(\lambda)$, arising as
the exponential rate of a restricted second moment, is negative except at
its two degenerate zeros $\lambda \in \{0, 1\}$. Huang's theorem assumes his
Condition~1.3: a two-variable function
$(\lambda_1,\lambda_2) \mapsto \mathcal{S}_\star(\lambda_1,\lambda_2)$,
the rate function of a planted first moment, is nonpositive on all of
$\R^2$, vanishing at the single point $(1,0)$. Both papers verified their
conditions numerically, in ordinary floating-point arithmetic, and stated
them as hypotheses. Huang's second arXiv version added rigorous interval
arithmetic for every numerical constant in his proof except Condition~1.3
itself; Ding and Sun wrote that they intended to supply a rigorous
verification in a later version of their paper, which did not appear.

This paper closes the remaining gap: we verify both conditions, and
re-derive the parameter rectangle
\[
  \alpha_\star \in [0.833078599,\, 0.833078600], \qquad
  q_\star \in [0.56394907949,\, 0.56394908030], \qquad
  \psi_\star \in [2.5763513100,\, 2.5763513224]
\]
of \cite[Proposition~1.3]{ding2018capacity}, on which every numerical
statement in both papers (and in this one) rests, in rigorous interval
arithmetic. Consequently:

\begin{theorem}\label{thm:main}
For the Ising perceptron at $\kappa = 0$ with Gaussian disorder,
$M_N/N \to \alpha_\star$ in probability. The same holds for i.i.d.
subgaussian disorder with mean zero and unit variance, in particular for
$\pm 1$ entries, by the universality theorem of \cite{nakajima2023sharp}.
\end{theorem}

The verification is computer-assisted in the sense that finitely many
inequalities between explicitly defined real numbers are established by
interval arithmetic (we use the ball arithmetic of Arb via python-flint,
with certified adaptive quadrature for all integrals). It is not, however,
a matter of running an integrator over the stated conditions. Three
genuinely mathematical obstructions have to be removed first, and most of
this paper is devoted to them.

First, Huang's condition is a statement over the whole plane, and the
supremum at infinity is not attained: $\mathcal{S}_\star$ tends to zero
along the degenerate directions only polynomially. Reparametrizing by the
moments of the pair $(X,\tanh X)$, $X\sim N(0,\psi_\star)$
(\S\ref{sec:huang-moments}), maps the plane onto a bounded convex body and
removes the tail entirely, while convex duality turns the entropy term into
a per-cell linear function; the bulk of the body is then certified by an
adaptive sweep (\S\ref{sec:huang-sweep}).

Second, both variational functions attain the value zero at degenerate
maximizers, where no direct interval evaluation can certify a sign, and near
Huang's maximizer no box enclosure of any Hessian of the problem can work
either: the tilted covariance $\nabla^2\Phi$ is nearly singular there, and
its inverse---which the curvature of the entropy term is---loses an
essential cancellation in interval arithmetic. We certify concavity of a
majorant along rays from the maximizer instead, bounding the entropy Hessian
by an explicit matrix over sublevel-set localizations of the dual variable;
nothing is ever inverted (\S\ref{sec:huang-sweep}). The pinned value and
gradient at the center are exact identities of the fixed point and require
no numerics. For the Ding--Sun condition the degenerate points are handled
by their own derivative and second-derivative analysis and a corrected
near-$\lambda=1$ chain, whose numerical inputs we certify (\S\ref{sec:ds},
\S\ref{sec:assembly}).

Third, the interval evaluations themselves have to be organized so that
widths do not defeat the margins, which are as small as $10^{-3}$ in the
interior of the Ding--Sun grid and near the star boundary of the Huang
sweep. Every delicate cell is evaluated in exact mean-value form---value at
the center plus an enclosed gradient times the radius---so that the
near-cancellations the margins depend on survive enclosure; the exact
identities $\mathcal{H}'(\lambda) = -(1-q_\star)\log A(\lambda)/2$ and the
explicit derivative integrals supply the gradients, and parameter-ball
evaluation subsumes the manual corner analyses mechanically.

\medskip
\noindent
The verification programs, the logs of the certified runs, and instructions
to reproduce every number in this paper are available at
\cite{repo}. All decimal constants
enter the programs as exact rationals; all comparisons are between balls
whose correctness is guaranteed by Arb. We follow the numbering of
\cite{ding2018capacity} and \cite{huang2024capacity} for their statements.

\section{The statements being verified}\label{sec:statements}

We use the normalizations of \cite{ding2018capacity,huang2024capacity}. With
$\varphi,\Psi$ the standard Gaussian density and upper tail, and
$\cE = \varphi/\Psi$, set for $q\in[0,1)$
\[
  F_{1-q}(x) = \frac{\cE}{\sqrt{1-q}}\Bigl(\frac{-x}{\sqrt{1-q}}\Bigr),
  \quad
  P(\psi) = \E\tanh^2(\sqrt\psi Z),
  \quad
  R_\alpha(q) = \alpha\,\E F_{1-q}(\sqrt q Z)^2,
\]
and the Gardner free energy $\sG(\alpha,q,\psi)$ of
\cite[eq.~(1.6)]{ding2018capacity}. The
threshold $\alpha_\star$ is the root of $\sG$ along the fixed point
$q=P(\psi),\ \psi=R_\alpha(q)$; Ding--Sun Proposition~1.3 makes this precise
and pins the rectangle of \S\ref{sec:gardner}. Huang's Condition~1.3 is
$\mathcal S_\star\le 0$, the object of \S\ref{sec:huang-moments}--%
\ref{sec:huang-sweep}; Ding--Sun's Condition~1.2 is
$\mathscr S_\star(\lambda)<0$ for $\lambda\notin\{0,1\}$, the object of
\S\ref{sec:ds}. \S\ref{sec:assembly} records which theorem consumes which
verified fact.

\section{The parameter rectangle}\label{sec:gardner}

Block~1 (\texttt{block1\_gardner.py}) re-proves Ding--Sun Proposition~1.3 in
interval arithmetic: thirteen certified inequalities establishing (i) that
$R_\alpha$ maps the stated $q$-intervals into the stated $\psi$-intervals;
(ii) the Almeida--Thouless contraction $\sup_{G}\,\mathrm dP(R_\alpha)/\mathrm dq
\le 0.96<1$, via $P'\le 0.08$ and $\mathrm dR/\mathrm dq\le 12$; (iii) the
fixed-point sign changes $P(R(q,\alpha))-q$ at the four corners; and (iv)
$\sG_\star(\alpha_{\mathrm{ub}})<0<\sG_\star(\alpha_{\mathrm{lb}})$. Each is a
one-dimensional certified Gaussian integral (adaptive \texttt{acb.integral}
plus an explicit tail bound), evaluated with $\alpha,q,\psi$ as balls covering
their whole ranges, which subsumes Ding--Sun's manual corner analysis. This
establishes $\alpha_\star\in[0.833078599,0.833078600]$ and the accompanying
$q_\star,\psi_\star$ intervals rigorously.

\section{The Ding--Sun condition}\label{sec:ds}

Ding--Sun's Condition~1.2, $\mathscr S_\star(\lambda)=\mathcal H(\lambda)
+\mathcal P(\lambda)+\mathcal A(\lambda)<0$ for $\lambda\notin\{0,1\}$, is
one-dimensional and mirrors Huang's: a value sweep on the bulk and a local
second-derivative analysis at the two degenerate zeros $\lambda=0,\lambda=1$.
Since $\mathcal A\le 0$, it suffices to bound $\mathcal H+\mathcal P$ on the
positive branch and, on the negative branch, the tilted majorant
$\mathcal H+\mathcal Q$ (Ding--Sun's $s=0.2$ shift). We build the certified
evaluators (\texttt{dsfun.py}): $\mathcal H(\lambda)=\mathfrak H(A(\lambda))$
through the pair-entropy $\Gamma$ and the map $A=\exp(2\,\mathrm{atanh}\,\lambda)$,
with the well-conditioned $D_H(A)=(1-m^2)\bigl(1-2/(\Delta+1)\bigr)$ that
removes the interval blow-up of the raw $(A^2-1)/(\Delta+1)^2$ form; and
$\mathcal P,\mathcal Q$ through the double integral $I_s(\lambda)$, evaluated
by nested certified quadrature with mean-value corrections. The bulk cells
verify. Near $\lambda=1$ we follow Ding--Sun's own chain (their Lemma~8.2,
Corollary~8.3, Propositions~8.4 and~8.1) with corrected constants: the two
integrals behind their Lemma~8.2 evaluate to $2.678\ldots$ and
$3.162\ldots$, not the printed $1.78$ and $4.3$ (the printed consequence
$1-\ell(A)\le 1.83/A$ is already inconsistent with their verified
$1-\ell(100)>0.025$), and the double integral of their Proposition~8.4 is
$-0.4447\ldots$ rather than $-0.45$. With the certified values
($1-\ell(A)\le 2.711/A$ for $A\ge100$, corollary constant $0.9987$, closing
coefficient $1.2837$) the same proof gives
$\mathcal H+\mathcal P<\mathcal H(1)+\mathcal P(1)$ on $[0.982,1)$, and the
part~(a) grid reaches $\lambda(0.99)=0.98665\ldots>0.982$, so the covering
of $[\lambda_{\min},1)$ is unchanged. The neighborhood of $\lambda=0$ is
handled by the derivative and second-derivative certificates described in
\S\ref{sec:assembly}.

\section{Huang's condition in moment coordinates}\label{sec:huang-moments}

Huang's Condition~1.3 asks that
\[
  \mathcal{S}_\star(\lambda_1,\lambda_2)
  = \inf_{s\ge 0}\Bigl\{ \tfrac12 s^2\psi_0 + \ent(\Lambda_{\lambda_1,\lambda_2})
    + \alpha_\star\, \E\log\Psi(V_s) \Bigr\} \le 0
  \qquad\text{for all }(\lambda_1,\lambda_2)\in\R^2,
\]
where $\Lambda_{\lambda_1,\lambda_2}(x)=\tanh(\lambda_1 x+\lambda_2\tanh x)$,
$X\sim N(0,\psi_0)$, $M=\tanh X$, and the constraint term depends on the
profile only through the two moments
\[
  a_1 = \E[X\,\Lambda], \qquad a_2 = \E[M\,\Lambda].
\]
The plane $\R^2$ of $(\lambda_1,\lambda_2)$ is unbounded, and the supremum is
not attained: along the degenerate directions $\mathcal{S}_\star\to 0^-$ only
polynomially. Verifying a nonpositivity statement over an unbounded domain by
interval arithmetic is the obstruction that stopped a direct check.

We remove it by changing coordinates. The profile
$\Lambda_{\lambda_1,\lambda_2}$ is exactly the maximum-entropy profile for its
own moments (this is the Lagrange condition that produced the two-parameter
family in the first place), so
\[
  \ent(\Lambda_{\lambda_1,\lambda_2}) = H(a_1,a_2)
  := \max\bigl\{\, \E\,\ent_2(\tfrac{1+\Lambda}2) :
     \E[X\Lambda]=a_1,\ \E[M\Lambda]=a_2,\ |\Lambda|\le 1 \,\bigr\},
\]
and therefore
\begin{equation}\label{eq:moment-reduction}
  \sup_{\lambda_1,\lambda_2} \mathcal{S}_\star(\lambda_1,\lambda_2)
  = \sup_{(a_1,a_2)\in K} \bigl[\, H(a_1,a_2) + G(a_1,a_2) \,\bigr],
  \qquad
  G(a_1,a_2) = \inf_{s\ge0}\bigl\{\tfrac12 s^2\psi_0
     + \alpha_\star\, \E\log\Psi(V_s)\bigr\}.
\end{equation}
The key point is that the new domain
\[
  K = \bigl\{ (\E[X\Lambda],\E[M\Lambda]) : |\Lambda|\le 1 \bigr\}
\]
is the \emph{moment body} of the pair $(X,M)$, a compact convex set with
support function $h(u,v)=\sup_{|\Lambda|\le1}\E[\Lambda(uX+vM)]=\E|uX+vM|$. It
is contained in the rectangle $|a_1|\le\E|X|$, $|a_2|\le\E|M|$. The unbounded
tail of the $(\lambda_1,\lambda_2)$ plane is compressed onto $\partial K$,
where the profiles are saturated ($|\Lambda|=1$) and hence
$H\!\restriction_{\partial K}=0$.

Two consequences make \eqref{eq:moment-reduction} checkable. First, the
constraint term $G$ depends on $(a_1,a_2)$ explicitly, through
$V_s = \bigl(-\tfrac{a_2}{q_0}\widehat H - \tfrac{a_1}{\psi_0}\bN\bigr)/D
 + s\bN$, $D=\sqrt{1-a_2^2/q_0}$, a one-dimensional Gaussian integral; and
$G\le \tfrac12 s^2\psi_0 + \alpha_\star\E\log\Psi(V_s)$ for \emph{any} fixed
$s$, so a per-cell choice of $s$ gives a certified upper bound with no inner
optimization. Second, convex duality gives, for \emph{every} dual point
$(b_1,b_2)$,
\begin{equation}\label{eq:dual-bound}
  H(a_1,a_2) \le \Phi(b_1,b_2) - b_1 a_1 - b_2 a_2,
  \qquad \Phi(b_1,b_2) = \E\log 2\cosh(b_1 X + b_2 M),
\end{equation}
with equality when $(b_1,b_2)$ is the dual of $(a_1,a_2)$. Choosing $(b_1,b_2)$
per cell (found by a Newton solve on the convex dual objective, then used as a
fixed rational) turns the entropy into an explicit linear function of
$(a_1,a_2)$, tight on the cell.

\begin{lemma}[Moment reduction]\label{lem:moment}
With $H$ and $G$ as above, $\sup_{\lambda_1,\lambda_2}\mathcal S_\star
=\sup_{(a_1,a_2)\in K}(H+G)$, and \eqref{eq:dual-bound} holds for all
$(b_1,b_2)$.
\end{lemma}
\begin{proof}
Fix $(\lambda_1,\lambda_2)$ and write $\Lambda=\Lambda_{\lambda_1,\lambda_2}$,
$a_i=a_i(\lambda)$. The constraint term of $\mathcal S_\star$ depends on
$\Lambda$ only through $(a_1,a_2)$ (it enters
$\mathcal S_\star(\Lambda,s)$ only via $\E[M\Lambda]=a_2$ and
$\E[X\Lambda]=a_1$, in the notation of \S\ref{sec:huang-moments}), so
$\mathcal S_\star(\lambda_1,\lambda_2)=\ent(\Lambda)+G(a_1,a_2)$. It remains
to show $\ent(\Lambda)=H(a_1,a_2)$, i.e. that $\Lambda$ maximizes the entropy
among all profiles with moments $(a_1,a_2)$. This is the Lagrange condition
that defines $\Lambda$: maximizing $\E\,\mathrm{ent}_2(\tfrac{1+\Lambda}2)$
over $|\Lambda|\le1$ subject to $\E[X\Lambda]=a_1,\E[M\Lambda]=a_2$ has,
by pointwise optimization of the Lagrangian
$\mathrm{ent}_2(\tfrac{1+\Lambda}2)+(\lambda_1 X+\lambda_2 M)\Lambda$, the
maximizer $\Lambda^\ast=\tanh(\lambda_1 X+\lambda_2 M)=\Lambda$. Hence
$\ent(\Lambda)=H(a_1,a_2)$ and
$\mathcal S_\star(\lambda_1,\lambda_2)=(H+G)(a_1(\lambda),a_2(\lambda))$.
Taking suprema and noting that $(a_1(\lambda),a_2(\lambda))$ ranges over a
dense subset of $K$ (every extreme point of $K$ is a limit of such profiles)
gives the first claim. For \eqref{eq:dual-bound}, strong duality for this
concave program reads $H(a)=\min_b[\Phi(b)-b\cdot a]$, where the inner
maximization of the Lagrangian over $|\Lambda|\le1$ gives
$\max_\Lambda[\mathrm{ent}_2(\tfrac{1+\Lambda}2)+\theta\Lambda]=\log2\cosh\theta$
at $\theta=b_1X+b_2M$; the stated inequality is the weak-duality direction,
valid for every $b$.
\end{proof}

Region~II certifies, cell by cell, the inequality
$\min_k[\Phi(b^k)-b^k\!\cdot a]+\tfrac12(s^C)^2\psi_0+\alpha_\star T_C(a)<0$ on
the cell, where $T_C$ encloses $\E\log\Psi(V_{s^C})$ over the cell by a
mean-value rule in $(a_1,a_2)$. By Lemma~\ref{lem:moment} the first term
bounds $H$ and the rest bounds $G$ from above on the cell, so the certified
inequality gives $(H+G)<0$ there. Cells whose closure misses $K$ are dropped:
a cell lies outside $K$ once some direction $(u,v)$ has
$\min_{a\in\mathrm{cell}}(u a_1+v a_2)>h(u,v)$, which is checked against a
certified upper bound on $h(u,v)=\E|uX+vM|$ over a fan of directions.

\section{The compact sweep and the degenerate point}\label{sec:huang-sweep}

Write $a^\star=(\psi_0(1-q_0),q_0)$ for the image of $(\lambda_1,\lambda_2)=(1,0)$;
it is the unique global maximizer, with $H(a^\star)+G(a^\star)=0$. We verify
\eqref{eq:moment-reduction} in two pieces.

\emph{Region II (the bulk).} For every cell $C$ of a bounded, adaptively
refined grid on $K$ away from the star region of Region~I, we certify
\[
  \bigl[\Phi(b^C) - b^C_1 a_1 - b^C_2 a_2\bigr]
  + \tfrac12 (s^C)^2\psi_0 + \alpha_\star\, T_C \;<\; 0
  \qquad\text{on } C,
\]
where $b^C,s^C$ are the (nonrigorous) per-cell dual and tilt and $T_C$ is a
mean-value enclosure of $\E\log\Psi(V_{s^C})$ over $C$. Cells lying outside
$K$ (detected by the support function $h$ along a fan of directions)
contribute nothing and are dropped; cells straddling $\partial K$ are handled
by pulling the dual inward, using that \eqref{eq:dual-bound} holds for any
fixed $b$. The sweep runs in two stages---all of $K$ minus a coarse box
around $a^\star$, then that box minus the star, bisecting to side $10^{-3}$
along the star boundary---and a cell is skipped only if a conservative test
places it entirely inside the star (maximal corner distance below the
minimal star radius over the cell's angular range). Both stages terminate
with every leaf cell strictly negative.

\emph{Region I (the maximizer).} On $E$ the value is genuinely $0$ at
$a^\star$, so no direct sign check can succeed there, and interval evaluation
of any Hessian of the problem over a two-dimensional box fails for a
structural reason: $\nabla^2\Phi(1,0)$ has condition number $\approx 82$
(the pair $(X,\tanh X)$ is strongly correlated), so the dual image of an
$a$-box is stretched by a factor $\approx 220$ along one direction and every
box enclosure of $\nabla^2 H=-[\nabla^2\Phi]^{-1}$ loses the cancellation it
needs. We work instead along rays. For a unit direction $v$ and a fixed
slope $\dot\sigma$ set
\[
  \varphi_v(t) = H(a^\star+tv) + \tfrac12 s(t)^2\psi_0
     + \alpha_\star T\bigl(a^\star+tv, s(t)\bigr),
  \qquad s(t)=s_0+t\dot\sigma,\quad s_0=\sqrt{1-q_0} .
\]
Since $G$ is a minimum over the tilt, $\varphi_v\ge \mathcal S_\star$ along
the ray, while at the fixed point (Huang's unconditional identities: the
Gardner stationarity and the tilt equation) $\varphi_v(0)=0$ and
$\varphi_v'(0)=0$. Hence certifying $\varphi_v''\le 0$ on $[0,T(v)]$ proves
$\mathcal S_\star\le0$ on the segment, and a fan of direction cells covering
the circle proves it on the star region
$\{a^\star+tv:\, 0\le t\le T(v)\}\supseteq E$. The second derivative splits
into a constraint part---certified fixed-grid integrals over the (thin) ray
piece, in which no dual variable appears---and the entropy part
$-v^\top[\nabla^2\Phi(\lambda(x))]^{-1}v$, which we bound WITHOUT tracking
the ill-conditioned dual map: since
$\nabla^2\Phi(\lambda)=\E[ff^\top\mathrm{sech}^2(\lambda\!\cdot\!f)]$ with
$f=(X,\tanh X)$, any localization $\lambda(x)\in\Lambda$ gives
\[
  \nabla^2\Phi(\lambda(x)) \preceq
  B_\Lambda := \E\bigl[\, f f^\top\, \max_{\lambda\in\Lambda}
     \mathrm{sech}^2(\lambda\cdot f)\bigr],
  \qquad\text{so}\qquad
  v^\top \nabla^2 H\, v \le -\, v^\top B_\Lambda^{-1} v ,
\]
and $B_\Lambda$ is a single explicit one-dimensional integral (the inner
maximum is a per-$z$ corner selection). With $\Lambda=\R^2$ this is the
global bound $B_0=\E ff^\top$, which settles every direction outside a
narrow window ($\approx\pm0.095$ rad) around the weak eigenvector of
$\nabla^2\mathcal S_\star(a^\star)$. Inside the window the localization
$\Lambda$ is supplied per radius band by a sublevel-set argument that never
inverts anything: $\lambda(x)$ minimizes $\Phi(\lambda)-\lambda\cdot x$, so
it lies in the convex set
$\{\lambda:\Phi(\lambda)-\lambda\cdot x\le \min_k[\Phi(\hat\lambda_k) -
\hat\lambda_k\cdot x]\}$ for any finite anchor fan $\hat\lambda_k$, and a
certified one-dimensional sweep of $\Phi$ over the boundary of a candidate
box (in exact mean-value form, with the per-segment derivative
$\nabla\Phi-x$ enclosed---it is genuinely small along the flat valley)
verifies that this set stays inside the box for every $x$ in the band; the
test function is convex in $x$, so corner points of the band hull suffice.
The boxes shrink quadratically-in-radius toward $(1,0)$, so
$B_\Lambda\to\nabla^2\Phi(1,0)$ exactly where sharpness is needed. All
certificates run over the parameter balls of the Ding--Sun rectangle, hence
cover the true fixed point.

\section{Assembly}\label{sec:assembly}

The pieces combine as follows.

\emph{Upper bound.} Huang's Theorem~\ref{thm:main} reduces the capacity upper
bound at $\kappa=0$ to four numerical conditions. His other three
numerical conditions (the parameter rectangle, the AMP constants, and the
local concavity constants) were verified by Huang in 	exttt{python-flint} and,
for the first, rest on Ding--Sun Proposition~1.3, which we re-verify in
\S\ref{sec:gardner} (Block~1). His Condition~1.3 is the statement
$\mathcal S_\star\le 0$ on $\R^2$. In moment coordinates
(\S\ref{sec:huang-moments}) this is $\sup_{K}(H+G)\le 0$; the two sweeps of
Region~II (\S\ref{sec:huang-sweep}) certify $H+G<0$ on $K$ outside the
certified star region around $a^\star$ (the first over $K$ minus a coarse
box, the second over that box minus the star, bisecting to side $10^{-3}$
along the star boundary), and Region~I's ray certificates give $H+G\le0$ on
the star itself. Hence $\sup_K(H+G)=0$, i.e.\
Condition~1.3 holds, and Huang's upper bound
$\lim_N \PP(M_N/N\ge\alpha)=0$ for $\alpha>\alpha_\star$ becomes
unconditional.

\emph{Lower bound.} Ding--Sun's main theorem gives, under their
Condition~1.2, that $\liminf_N\PP(M_N/N\ge\alpha)>0$ for $\alpha<\alpha_\star$;
the sharp-threshold results of Xu and Nakajima--Sun upgrade this to high
probability. Condition~1.2 is the one-dimensional statement
$\mathscr S_\star(\lambda)<0$ for $\lambda\notin\{0,1\}$, and the covering of
$[\lambda_{\min},1]$ follows their own decomposition: the value grids on
$[0.2,0.98]$ and $[\lambda_{\min},-0.125]$ (part (a), our Block~3a), the
derivative signs $\mathrm{PG}'<0$ on $[0.05,0.2]$ and $\mathrm{PG}'>0$ on
$[-0.125,-0.03]$ (part (b)), the second derivative $\mathrm{PG}''<0$ on
$[-0.03,0.05]$ around the degenerate zero $\lambda=0$ where
$\mathrm{PG}(0)=\mathrm{PG}'(0)=0$ (part (c)), and the near-one analysis on
$[0.98,1)$ (Block~2). Parts (b) and (c) use the exact derivative identities
$\mathcal H'=-(1-q_\star)\log A/2$, $\mathcal
H''=-(1-q_\star)/(2A\,\ell'(A))$ and certified two-dimensional integrals for
$I'$ and $I''$ (with closed-form Gaussian-moment bounds for the mass outside
the integration box); the interval endpoints are pinned by the certified
monotone map $\lambda=\ell(A(\tau))$, exactly as in part (a).

\emph{Disorder.} Krauth and M\'ezard's model has $\pm1$ (Bernoulli) disorder;
we work with Gaussian disorder throughout. The two agree because the capacity
has a universal sharp threshold sequence over subgaussian disorder
(Nakajima--Sun), which includes both.

\emph{What is machine-checked.} Every statement labelled a Claim is a finite
list of ball inequalities between explicitly defined real numbers, established
by Arb; the analytic reductions around them (the moment-body reduction, convex
duality, the ray majorants and their pinned identities, the localization and
$B_\Lambda$ arguments, the reductions to certified quadrature) are proved on
paper. The nonrigorous companion (\texttt{huang\_np.py}) only selects
per-cell duals, tilts, anchors, and candidate boxes, which enter the
certificates as fixed rationals; no certified inequality depends on it.

\emph{Inventory.} The final runs comprise: the parameter rectangle,
13~checks; Huang Region~II, $873$ leaf cells (stage~1, the moment body minus
a coarse box) and $1436$ leaf cells (stage~2, the box minus the star), all
strictly negative; Huang Region~I, $836$ band certificates covering the star
with zero failures; Ding--Sun part~(a), $118$ grid cells; the corrected
near-one block, $10$ checks; and the middle-interval derivative and
second-derivative certificates. Logs and JSON summaries for every run are
included with the source.

\begin{remark}[Soundness audit]
Two implementation errors were caught during this work by independent
cross-checks and corrected before the final runs; we record them because
they illustrate what interval arithmetic does and does not protect against.
First, an early version of the fixed-grid quadrature multiplied the
derivative enclosure of the mean-value remainder by the \emph{signed} first
moment $\int_{\text{cell}}(z-m)\varphi\,dz\approx 0$; since the enclosed
derivative varies over the cell while $(z-m)$ changes sign, the remainder
does not factor through the signed moment, and the resulting balls were
systematically too tight---detected because a 30-digit quadrature of
$\nabla\Phi(1,0)$ fell outside them. The corrected rule handles the
positive and negative parts of $(z-m)$ separately. Second, a sign error in
$\partial_s^2$ of the tilted constraint term ($+$ for $-$) survived until a
finite-difference cross-check of every ingredient of $\varphi_v''$ against
an independent floating-point implementation. Neither error is visible to
the interval machinery itself: Arb guarantees the arithmetic, not the
formulas. All certificates reported here postdate both corrections.
\end{remark}

\section*{Funding and declarations}
The funding and grant information should be verified by the author before
submission. In prior work the author has listed support from the European
Research Council under the European Union's Horizon 2022 research and
innovation programme, the Simons Foundation, Heights Labs, and Israel Science
Foundation grants 2258/19 and 4101/25.

The large majority of the work reported here---the mathematical strategy,
the design of the interval certificates, essentially all of the code, the
soundness audits, and the drafting of this manuscript---was carried out by
Fable~5 (Anthropic), an AI system, operating in an autonomous goal mode: the
author set the objective of resolving the outstanding conditions, and the
system worked toward it over extended unattended sessions, choosing the
mathematical approach, implementing and running the certificates, detecting
and repairing its own errors through independent cross-checks, and iterating
until the verifications closed. The author supervised the runs, reviewed the
mathematics, and accepts responsibility for the final manuscript. This is
part of a broader research project of the author on the extent to which
current AI systems can resolve open problems in mathematics.

\begin{thebibliography}{9}

\bibitem{krauth1989storage}
W.~Krauth and M.~M\'ezard,
Storage capacity of memory networks with binary couplings,
J. Physique \textbf{50} (1989), 3057--3066.

\bibitem{ding2018capacity}
J.~Ding and N.~Sun,
Capacity lower bound for the Ising perceptron,
STOC 2019; arXiv:1809.07742.

\bibitem{xu2021sharp}
C.~Xu,
Sharp threshold for the Ising perceptron model,
Ann. Probab. \textbf{49} (2021), 2399--2415.

\bibitem{nakajima2023sharp}
S.~Nakajima and N.~Sun,
Sharp threshold sequence and universality for Ising perceptron models,
SODA 2023; arXiv:2204.03469.

\bibitem{huang2024capacity}
B.~Huang,
Capacity threshold for the Ising perceptron,
FOCS 2024; arXiv:2404.18902.

\bibitem{altschuler2024capacity}
D.~J.~Altschuler and K.~Tikhomirov,
A note on the capacity of the binary perceptron,
arXiv:2401.15092.

\bibitem{repo}
Y.~Shmalo,
Verification programs and certified run logs for this paper,
\url{https://github.com/yspennstate/ising-perceptron-capacity}, 2026.

\end{thebibliography}

\end{document}
